%======================================================================
\section{Custo do sistema}\label{sec:custo}
%======================================================================
Cada quarto recebe um hub e uma pulseira. A Tabela~\ref{tab:custo_resumo} resume quanto custa
equipá-lo, por grupo de itens, e a Tabela~\ref{tab:custo}, no Apêndice~\ref{ap:bom}, traz
cada item com a loja e o preço.

\begin{table}[H]
\centering
\caption{Custo estimado de um quarto, por grupo de itens.}
\label{tab:custo_resumo}
\small
% Gerado por modelo/experimentos/12_relatorio_q2.py a partir de hardware/placas/v1/custo-estimado.csv
\begin{tabular}{@{}l r r r r@{}}
\toprule
Grupo & Hub (R\$) & Pulseira (R\$) & Kit (R\$) & Parte do kit \\
\midrule
ESP32-S3-DevKitC-1 (uma por placa) & 94,30 & 94,30 & 188,60 & 33\,\% \\
PCB v1 (estimada) & 35,00 & 35,00 & 70,00 & 12\,\% \\
Módulos, sensores das portas e cabo & 60,42 & 22,80 & 83,22 & 15\,\% \\
Componentes da placa & 17,86 & 2,87 & 20,73 & 4\,\% \\
Alimentação & 34,43 & 73,04 & 107,47 & 19\,\% \\
Caixas & 46,91 & 46,91 & 93,82 & 17\,\% \\
\midrule
\textbf{Total} & \textbf{288,92} & \textbf{274,92} & \textbf{563,84} & 100\,\% \\
\bottomrule
\end{tabular}

\end{table}

\paragraph{Critério.} Os preços são do varejo brasileiro, para compra por unidade, sem frete,
consultados nos sites das lojas em \custoData. Cada PCB leva só as peças da sua função: o hub,
o LM35, o LM358, o PIR, o CAN e as portas; a pulseira, o MPU e o botão. A PCB entra como
estimativa, a partir do preço da JLCPCB com frete e impostos de importação
(Apêndice~\ref{ap:bom}). Ficam de fora a central, um único Raspberry Pi para todos os quartos
da ILPI, e a montagem, feita à mão.

\paragraph{O que pesa.} Um quarto custa cerca de R\$\,\custoKit: R\$\,\custoHub{} o hub e
R\$\,\custoPulseira{} a pulseira. As duas DevKits são \custoDevkitPct\,\% desse valor, e a
alimentação e as caixas, outros \custoAlimCaixasPct\,\%. Os sensores, os módulos e os
componentes da placa, que fazem o trabalho do sistema, somam menos de um quinto. A maior parte
do custo do protótipo está nas peças prontas, compradas para encurtar o desenvolvimento, e não
nos circuitos projetados.

\paragraph{Em série.} Esse é o custo de um protótipo comprado no varejo, e ele cairia por três
caminhos:
\begin{itemize}
  \item trocar a DevKit pelo módulo ESP32-S3 soldado numa placa própria, que dispensa a placa de
        desenvolvimento inteira (regulador, conversor USB--serial, conectores);
  \item comprar módulos e componentes em volume, de distribuidores, em vez do varejo;
  \item dar à pulseira uma placa pequena, com bateria e carregador, no lugar do
        \textit{power bank} e da caixa do tamanho da do hub (Seção~\ref{sec:conclusoes}).
\end{itemize}
A caixa do hub também precisa de vedação nas aberturas, do PIR e da USB, para chegar ao IP42
pedido no RNF12.
